\section{Limitations}
\label{sec:limits}

\textbf{Software safety is not certified.} The protective field, watchdogs and e-stop run in
Python over DDS. They are a second channel; a robot in a public space needs a first channel
that is safety-rated (hardware e-stop and a certified safety scanner or bumper wired to the
drive's safe torque off), designed and validated to ISO 13849-1~\cite{iso13849} or IEC
61508~\cite{iec61508}, and a safety case (hazards, mitigations, verification). Dharma reads the
hardware e-stop state so both channels latch together, but it does not replace it.

\textbf{Perception gaps.} A Mid-360 sees down to $-7^\circ$: mounted 0.53~m above the floor, a
10~cm object is seen only beyond about 3.5~m. The obstacle memory keeps objects seen before, but
an object that appears inside this near-field zone is not detected, and missing floor (stairs
down, platform edges) is not an obstacle at all. A pitched mount or a near-field sensor is
required where this matters. The scan is not deskewed (up to $v \cdot 0.1$~m of distortion per
frame), which the field margins must cover.

\textbf{No loop closure.} Mapping integrates scans at the corrected odometry pose; large sites
accumulate error between fiducials. A pose graph over scan matches with tags as landmarks would
unify mapping and localization.

\textbf{Sim-to-real gap.} The evaluation is in simulation with modelled drift, latency and
dropout. Braking deceleration and reaction time must be measured on each robot (worst-case
floor, payload and battery) before the field limits can be trusted; the results do not
transfer to another robot without repeating them, and real-robot trials are future work.
